% 第 43 章　量子态不是一颗偷偷走固定路线的小球
\chapter{量子态不是一颗偷偷走固定路线的小球}

第 42 章的结尾留下了一桩悬案：电子——这颗在阴极射线管里被电场推着走、被当成小颗粒三十年的东西——穿过晶体后会衍射。会衍射是波的签名。可同一颗电子打在荧光屏上，永远是一个完整的点，从没有人见过“半颗电子”。这一章，我们把镜头推进到那台把所有人都逼疯的实验装置：双缝。它只用两条狭缝和一块屏幕，就把经典物理最深的一个默认逼到了墙角——那个默认是：任何物体在任何时候都走在一条确定的路线上，区别只在于我们看没看见。

本章的任务是建立一套新的记账语言。这套语言只有两条规则，简单得可以写在一行里；但它的代价是，我们必须发明一个比“概率”更原始的概念——概率幅——并且承认：量子态不是一颗偷偷走固定路线的小球。

\step{反常识开场}

把一台电子显微镜改造成这样：电子源调到极弱极弱，弱到平均每秒只放出几千颗电子，而电子飞过整个装置只需要不到一亿分之一秒——算下来，任意时刻装置里平均连一颗电子都没有。两颗电子之间隔着漫长的时间荒漠，它们不可能相遇、碰撞、商量。

电子源前面立一块挡板，板上刻两条极细的平行狭缝；挡板后面放一块荧光探测屏，电子打上去会闪一下，并且仪器能记下每个闪光点的位置。现在，开始放电子，一颗，又一颗。

\begin{tikzpicture}[scale=0.95,>={Stealth}]
% 电子源
\fill (0,0) circle (1.5pt) node[left=2pt]{\small 电子源};
% 双缝挡板
\draw[thick] (3,-2) -- (3,-0.62);
\draw[thick] (3,-0.38) -- (3,0.38);
\draw[thick] (3,0.62) -- (3,2);
\node at (3.05,1.6) [right]{\small 双缝};
% 两条虚线路径
\draw[dashed] (0,0) -- (3,0.5) -- (7,1.05);
\draw[dashed] (0,0) -- (3,-0.5) -- (7,1.05);
% 探测屏
\draw[thick] (7,-2) -- (7,2);
\node at (7,2.15) {\small 探测屏};
% 屏上条纹（长度示意亮度）
\foreach \y/\w in {-1.5/0.22,-0.75/0.62,0/0.95,0.75/0.62,1.5/0.22}
  \fill[gray!60] (7.12,\y-0.07) rectangle (7.12+\w,\y+0.07);
\node at (7.9,-1.9) [right]{\small 条纹};
\end{tikzpicture}

第一颗电子：屏上闪了一个点，位置看起来毫无道理。第二颗：另一个点，离得挺远。前几十颗落点像随手撒下的一把盐，看不出任何规律。几百颗之后，隐约有点不对：某些区域点子越来越密，某些区域始终空空如也。几万颗之后，照片上分明的结构浮出水面——明暗相间的条纹，一条亮，一条暗，整齐地排列下去，和光的双缝干涉图样一模一样。

现在盯住一条暗带。那是几万颗电子里几乎没有一颗落下的地方。做两个对照实验。第一个：把右边的缝挡住，只留左边的缝，重新放电子——暗带那个位置，有电子落下了，那里属于单缝图样的亮区。第二个：换过来，只留右边的缝——那个位置同样有电子落下。两个缝单独开着时都能到达的地方，两条缝一起开时，反而到不了了。

多开一条路，到达反而变少。

再强调一次：电子是一颗一颗来的，每颗之间可以隔上整整一分钟，实验做上几个月，条纹照旧出现。没有任何“别的电子”可以替它探路、给它让路、和它干涉。每一颗电子都是独自面对那两条缝的——而它落点的规律，清清楚楚地表明它“知道”两条缝都开着。

费曼在《物理学讲义》里谈到这个实验时说过一句常被引用的话：双缝实验里藏着量子力学的“心脏”，它是量子力学唯一的谜——因为量子世界里其他所有的怪事，翻来覆去追到最后，都会回到这同一个谜上。这句话不是修辞，是经验总结。

\step{先猜一猜}

照例，拿纸写下你的答案，本章结尾回来打分。

第一题：电子一颗一颗地发，积攒几十万颗之后，你猜屏幕上是两条亮带（对应两条缝），还是一片均匀的小山包，还是明暗相间的多条条纹？

第二题：在暗带的某个位置，几千颗电子里几乎没有一颗落在那里。现在挡住其中一条缝，你猜那个位置的计数会（a）依然为零，（b）变多，（c）变多但不超过原来的一半？

第三题，也是本章真正的问题：对于那一颗一颗穿过装置的电子，你倾向于下面哪种图像？（a）每颗电子其实只走了某一条缝，只是我们没有看见是哪条；（b）电子像一小片云那样散开，一半走左缝一半走右缝，到屏上再合起来；（c）两种都不对。

写好了吗？如果你选了（a），你和几乎所有人——包括年轻时的爱因斯坦——站在一起。而这个看起来最稳妥、最讲理的答案，恰恰是错得最彻底的一个。

\step{动手实验}

家里做不了电子双缝，但“两个源头叠加、某些地方反而安静”这件事，两只音箱就能演给你看。

\begin{experiment}[两只音箱的静音点]
\textbf{材料：}两部手机（或一部手机加一只蓝牙小音箱），一个能发出稳定单音的手机应用或网页（搜索“在线正弦波发生器”即可），把频率设为一千赫兹。

\textbf{做法：}让两个声源播放同一个音，并排放在桌上，相距一米左右，音量一致。捂住一只耳朵，另一只耳朵在声源前方一两米的范围内缓慢移动，头部左右平移，仔细听音量的起伏。你会找到一串“响—轻—响—轻”交替的位置，其中某些点声音明显塌下去，几乎听不见——这就是静音点。一千赫兹的声音波长约为 \(34\) 厘米，相邻静音点的间距在这个量级，房间大小正合适。

\textbf{关键一步：}把头固定在静音点上，请人用手或一本书严严实实挡住其中一只音箱。你猜会怎样？那个几乎听不见的声音，会重新变响。撤掉遮挡，它又沉下去。多开一个声源，反而安静；关掉一个声源，反而热闹——和开场那个暗带变亮的故事，同一个剧本。

\textbf{为什么：}声波是空气压强的起伏，两列波在同一地点叠加时，压强的变化量相加。如果一列波恰好把空气“挤”的时候另一列在“拉”，两相抵消，压强几乎不动，耳朵就听不到什么。抵消是否发生，取决于两列波到达该处时“节拍”差多少——这就是相位差，它由路程差决定：路程差恰好是波长的整数倍，两列波同步，格外响；差半个波长，两列波对着干，格外安静。

\textbf{想一想：}静音点上“消失的声音”去了哪里？没有消失——亮的地方比一只音箱单独工作时更亮，能量从暗处搬到了亮处，总账不变。请记住这条守恒，本章后面要原样用到电子身上。
\end{experiment}

这个实验和电子双缝像在哪里：都是“两路贡献相加，同相增强、反相抵消”。不像在哪里：声波的振幅是可以拿麦克风直接测出来的空气压强起伏，是真真切切的物理量；而本章后面要出场的电子“振幅”，你拿任何仪器都测不到它本身。这个差别和相似处同样重要，请把它带在身边往下读。

\step{旧模型遇到麻烦}

先把实验事实钉在墙上，一共三条，任何模型都不许和它们讨价还价：

\begin{itemize}[leftmargin=1.8em]
\item 探测总是整颗的：每次闪光是一个点，每颗电子的电荷、质量、能量永远是完整的一份，从没见过半颗电子同时出现在两处；
\item 积累图样是条纹：大量电子的落点分布出现明暗相间的干涉结构；
\item 与数量无关：电子一颗一颗来、间隔再久，条纹照旧——每颗电子各自独立完成这件事。
\end{itemize}

现在让三个候选模型依次上场受审。

\textbf{模型一：电子是颗小球，每次走一条缝。}那么屏幕后方每个位置，电子只能从左缝或右缝到达，总的落点分布应该是“走左缝的分布”加“走右缝的分布”——两团小山包，充其量中间交叠。条纹，尤其是“两条缝都开时反而到不了”的暗带，无论如何出不来。模型一不及格。想打补丁“电子在半路上互相碰撞、挤出了条纹”？逐个发送的实验把这个补丁撕掉了：装置里从头到尾只有一颗电子，它和谁碰撞？

\textbf{模型二：电子是一小片波，裂成两半，各走一条缝。}这个模型能解释条纹——两列波在屏前叠加干涉。但它立刻撞上事实一：如果电子真的裂成两半，那么至少偶尔，屏上应该两处同时闪、每处带半颗电子的电荷；或者把两个探测器分别放在两条缝后面，它们应该常常同时轻响各记半份。实验的回答斩钉截铁：从来没有。探测永远是一整个点、一整份电荷、一声完整的响。电子在“被探测”这件事上从不分裂。模型二也不及格。

\textbf{模型三：电子每次只走一条缝，只是我们没有看见——走左走右是个概率问题。}这是几乎所有人的直觉答案，也是错得最脆生的一个，值得单独审判。想象一枚硬币扣在掌心：是正面还是反面，你“不知道”，但世界上存在一个确定的答案，你的概率只是无知的记账。对这种情况，概率论有一条铁律：到达某处的总概率，等于“走左缝到达”的概率加“走右缝到达”的概率，
\[
P_{12}(x) = P_1(x) + P_2(x).
\]
这条规则没有第二种可能——“走左”和“走右”穷尽且互斥，概率当然相加。可是实验说：在暗带处 \(P_1 > 0\)，\(P_2 > 0\)，而 \(P_{12} = 0\)。两个正数相加等于零，概率加法被判死刑。如果连“它到底走了哪条缝，只是我们不知道”都救不了场，那就只剩一条路：放弃“电子每时每刻都在走某一条确定路线”这个信念本身。

三个模型全部阵亡，但阵亡的方式指明了生路。模型三的失败告诉我们：需要的量必须能“抵消”，而概率永远非负，抵消不了。模型二的失败告诉我们：这个量不能是可以直接测到的实物——电子本身不散开。模型一的失败告诉我们：这个新量必须对“每一条可能的路径”都有响应，哪怕那条路径上没有“真的”有东西在走。

\step{发明新概念}

历史上，这套新语言是费曼整理得最干净的。全部规则只有三条：

\textbf{规则一：每一种“原则上不可区分”的发生方式，配一个箭头。}电子到达屏上某点，可以“经左缝到达”，也可以“经右缝到达”。如果没有任何装置、任何记录能把这两种方式区分开，那就给每种方式画一支箭头。箭头有两个属性：长度和指向。这个箭头叫概率幅；它的指向有个专门的名字，叫相位。

\textbf{规则二：不可区分的方式，箭头相加。}把两支箭头首尾相接，从起点画到终点，得到合成箭头——就是平面上的向量加法，和你画位移合成时用的方法一模一样。相位在这里开始发威：两支箭头指向相同时，合成的最长；指向相反时，合成缩成一个点；指向成一个角度时，介于两者之间。这就是叠加原理：量子态的“加”是箭头加，不是数字加。

\textbf{规则三：概率等于合成箭头长度的平方。}这条叫玻恩规则。注意是长度的平方——箭头指向哪个方向，对概率毫无影响；指向只在与别的箭头相加时才显露意义。可观测的世界（屏幕上闪不闪光）只认识箭头的长度，相位本身永远躲在幕后，只通过“相加”这个环节间接现身。

把这包规则推广开，就得到了本章的主角。问“电子可能在屏上哪些位置被探测到”，就对一个一个位置 \(x\) 各配一支箭头，记作 \(\psi(x)\)，读作波函数；\(\psi(x)\) 的长度平方给出在 \(x\) 附近探测到电子的概率（严格说是概率密度）。屏上一点、缝里一点、半路上一点，每处都有一支箭头，合起来构成这整套装置里电子的量子态——量子态不是电子的“藏身之处清单”，而是对“关于这个系统的任何测量，能测到什么、各有多大概率”的完整回答。

在这里，必须郑重地把三样东西分开，它们经常被通俗读物煮成一锅粥。

\textbf{实验事实}：屏上逐个的点、积累出的条纹、暗带关缝变亮。这些不依赖任何理论，换谁来重做都一样。

\textbf{数学模型}：箭头、相加规则、模的平方、波函数 \(\psi(x)\)。这套东西是人为发明的计算规则，它的全部资格来自“算出来的概率和实验逐点吻合”。

\textbf{物理解释}：箭头到底是什么？波函数是真实存在的东西，还是只是我们的记账工具？一颗电子“在测量之前”有没有位置？注意，以上三条规则一个字的解释都没给——它们只告诉你怎么算，不告诉你“世界内部在演什么戏”。物理学家为这出戏吵了一百年，提出过好几套互不相同的剧本，它们算出的概率一模一样。这幅诠释地图是第 45 章的内容；本章只需要记住一条纪律：别把任何一套剧本当成实验事实来宣讲。

顺便回头看看全书反复追问的三个问题，量子力学对第一个问题给出了崭新的答案。“系统此刻是什么状态？”——在力学里，答案是一组位置和动量；在热学里，是一组宏观参量；而在这里，答案是一个量子态：一包足以算出一切测量结果概率的振幅。“什么规律决定状态如何变化？”——留给第 44 章的薛定谔方程。“我们怎样通过测量知道它？”——就是玻恩规则，以及第 45 章要细讲的测量过程。三个老问题一个都没丢，只是答案的语法变了。

最后给概率幅一个声明完整的比喻。它像水波的波高：可以叠加，可以一峰一谷相互抵消，干涉条纹就是这么来的。它不像水波的波高：波高是可以拿尺子量的实物起伏，而概率幅测不到，它可以是负的，甚至（马上就会看到）需要两个数才能写全；问“这个位置的振幅本身是多大”没有操作上的意义，你能数的永远只有一颗颗电子的到达次数。它是计算概率的中转站，不是概率本身，更不是缩了水的物质波。

\step{一句话模型}

\begin{oneline}
每一种原则上分不出的路径贡献一支带相位的箭头，箭头先相加、再取长度的平方，才得到概率；所以量子态不是一颗偷偷走固定路线的小球——“走了哪条路”这个问题，在没有记录的世界里根本没有答案。
\end{oneline}

\step{数学翻译器}

把三条规则写成式子，并且每条都当场用极端情形检查。

\textbf{第一式：振幅相加，概率平方。}设两条路径的振幅为 \(A_1\) 和 \(A_2\)，则
\[
A = A_1 + A_2, \qquad P = |A|^2.
\]
竖线是箭头的长度。拿等长的两支箭头（长度都是 \(a\)）做三个特殊检查——同相：合成长度 \(2a\)，概率 \(4a^2\)。注意，是每条路径单独概率 \(a^2\) 的四倍，而不是两倍！“两条缝都开”在亮纹中心比“两条缝各自的贡献相加”还要亮一倍，这正是多出来的那一倍亮，对应着暗带被搬走的那份。反相：两支箭头头顶头，合成长度为零，\(P = 0\)——两个“都可能到达”加出“必然到不了”，概率加法做不到的事，振幅加法一笔完成。垂直：由勾股定理，合成长度 \(\sqrt{2}\,a\)，概率 \(2a^2\)，恰好等于两条路径概率之和。看出门道了吗？经典的概率相加 \(P = P_1 + P_2\) 并不是错，它只是相位差恰好九十度时的特例——它不是公理，而是更一般的规则在某个角度上的投影。

\begin{tikzpicture}[>={Stealth},scale=1.05]
% 同相
\draw[->,thick] (0,0) -- (1,0);
\draw[->,thick] (1,0) -- (2,0);
\draw[->,very thick] (0,-0.45) -- (2,-0.45);
\node at (1,-0.95) {\small 同相：长 \(2a\)，\(P=4a^2\)};
% 反相
\draw[->,thick] (4,0) -- (5,0);
\draw[->,thick] (5,0) -- (4,0);
\fill (4,-0.45) circle (1.2pt);
\node at (4.5,-0.95) {\small 反相：长 \(0\)，\(P=0\)};
% 垂直
\draw[->,thick] (8,0) -- (9,0);
\draw[->,thick] (9,0) -- (9,1);
\draw[->,very thick,dashed] (8,0) -- (9,1);
\node at (8.75,-0.95) {\small 垂直：长 \(\sqrt{2}\,a\)，\(P=2a^2\)};
\end{tikzpicture}

\textbf{第二式：条纹在哪里。}相位差从哪里来？在双缝这类实验里，主要来自路程差。第 42 章已经给出粒子的波长 \(\lambda = h/p\)；一路比另一路多走 \(\Delta s\)，相位就差
\[
\Delta\varphi = \frac{2\pi\,\Delta s}{\lambda}.
\]
同相（\(\Delta s = 0, \lambda, 2\lambda, \dots\)）处最亮，反相（\(\Delta s = \lambda/2, 3\lambda/2, \dots\)）处全暗。从几何关系可以推出屏上相邻亮纹的间距
\[
\Delta y = \frac{\lambda L}{d},
\]
其中 \(d\) 是双缝间距，\(L\) 是挡板到屏的距离。检查特殊情况——\(L\) 或 \(\lambda\) 翻倍：条纹间距翻倍，波越长、屏越远，图样摊得越开，合理；\(d\) 翻倍：间距减半，缝离得越远条纹越挤，合理；\(\lambda \to 0\)：\(\Delta y \to 0\)，条纹密到任何探测器都分辨不开，看起来就是一团连续的小山包——记住这个极限，它后面负责解释“你为什么看不见棒球的干涉”。

\textbf{一个带真实数据的估算：条纹到底有多宽。}取一百伏电压加速的电子，第 42 章算过：动量约 \(5.4\times10^{-24}\,\mathrm{kg\cdot m/s}\)，波长 \(\lambda \approx 0.12\,\mathrm{nm}\)。取双缝间距 \(d = 100\,\mathrm{nm}\)（大约是病毒的大小），屏放在 \(L = 10\,\mathrm{cm}\) 外，则
\[
\Delta y = \frac{1.2\times10^{-10}\times0.1}{1.0\times10^{-7}} \approx 1.2\times10^{-4}\,\mathrm{m}.
\]
十分之一毫米量级。所以电子双缝的条纹不会肉眼可见地铺在墙上，但在电子显微镜的放大下清清楚楚——开场那张“几万颗电子织出条纹”的照片，就是这么拍出来的，一九八九年前后由外村彰的团队做成了教科书级的版本。这个估算也告诉你实验难在哪：\(d\) 必须是纳米级，缝本身还得更窄，任何一个尺寸放大一百倍，条纹就密到糊成一片。

\begin{tikzpicture}[scale=0.92]
% 左：概率相加
\draw[->] (-0.2,0) -- (6.3,0) node[right]{\small 位置};
\draw[->] (0,-0.2) -- (0,1.4) node[above]{\small \(P\)};
\draw[domain=0.3:5.7,samples=60,smooth] plot (\x,{exp(-((\x-2.0)^2)/0.5)+exp(-((\x-4.0)^2)/0.5)});
\node at (3,-0.85) {\small 若概率相加：只有两团，无暗带};
% 右：振幅相加再平方
\begin{scope}[xshift=7.6cm]
\draw[->] (-0.2,0) -- (6.3,0) node[right]{\small 位置};
\draw[->] (0,-0.2) -- (0,1.4) node[above]{\small \(P\)};
\draw[domain=0.3:5.7,samples=120,smooth] plot (\x,{0.5*(1+cos(6*\x r))*(0.35+0.65*exp(-((\x-3.0)^2)/4))});
\node at (3,-0.85) {\small 振幅相加再平方：明暗相间};
\end{scope}
\end{tikzpicture}

\textbf{第三式：箭头为什么是两个数——复数的登场。}一支平面箭头，你总要报两个数才能讲清楚：横着走多少，竖着走多少。数学家早就给“一对数加上箭头加法”起好了名字：复数。把横分量记作 \(a\)，竖分量记作 \(b\)，箭头就写作 \(a + bi\)，其中 \(i\) 的定义只有一句话：乘一次 \(i\) 等于逆时针转九十度。转两次，\(i \times i = -1\)，方向整个反过来——这就是那个著名的 \(i^2 = -1\)，它不是玄学，是“转两个九十度”的记账。为什么这套语言对相位最自然？因为相位的本质就是转动，而复数的乘法规则自动替你处理转动的叠加：描述“先转这么多、再转那么多”，两个复数一乘就完事，不用背任何新的口诀。如果只用实数记账，你手里只有一个负号，只能表达“同向”和“反向”两档，九十度、三十七度这些连续变化的相位根本无处安放。所以复数不是物理学家故意附庸风雅，而是“箭头加转动”这件事的母语。顺便立一块护栏：概率幅里的“虚数”不代表它“虚幻”，正如负数出现在账目里不代表欠债是假的；可测量的永远是模的平方，是实得不能再实的计数。

\begin{mathbridge}
\textbf{把平方展开，看看“干涉项”长什么样。}设两路振幅大小为 \(|A_1| = a_1\)、\(|A_2| = a_2\)，相位差为 \(\Delta\varphi\)。用向量求模长的公式（或余弦定理），
\[
P = |A_1 + A_2|^2 = a_1^2 + a_2^2 + 2a_1a_2\cos\Delta\varphi
   = P_1 + P_2 + 2\sqrt{P_1P_2}\cos\Delta\varphi.
\]
前两项正是经典的概率相加；第三项叫干涉项，是量子世界多出来的那一栏账。逐项检查——\(\Delta\varphi = 0\)：\(\cos = 1\)，\(P = (\sqrt{P_1}+\sqrt{P_2})^2\)，相长，亮纹中心；\(\Delta\varphi = \pi\)：\(\cos = -1\)，\(P = (\sqrt{P_1}-\sqrt{P_2})^2\)，两路等强时恰好为零，暗带；\(\Delta\varphi = \pi/2\)：干涉项消失，回到经典相加。最有意思的是第四种情形：如果相位差因为某种原因变得杂乱无章、时快时慢，那么对时间或大量粒子平均后，\(\cos\Delta\varphi\) 正负相抵平均为零，干涉项自动蒸发，概率加法复活。请记住这个“平均为零”——第 45 章讲测量时你会看到，宏观世界之所以显得经典，很大程度上就是因为相位被环境洗乱了。

复数记法下，振幅写成 \(A = a + bi\)，长度平方 \(|A|^2 = a^2 + b^2\)；若只关心“长度为一、相位为 \(\theta\)”的箭头，写成 \(A = \cos\theta + i\sin\theta\)。波函数的归一化条件是
\[
\int |\psi(x)|^2 \, dx = 1,
\]
意思是“电子总归要在某处被探测到”，所有位置的概率密度加起来等于一。这是一条记账纪律：粒子数不会凭空多出来。
\end{mathbridge}

\begin{challenge}
\textbf{为什么不能只用概率，以及箭头随时间怎么转。}

假设有人不服气，说“我偏要只用概率描述双缝”。他的理论里，每个位置只有一个非负的数，不同路径的贡献只能相加——非负数加非负数永远还是非负数，数学上根本写不出“抵消”二字。要让“抵消”可能，贡献量必须能取负；而要让“抵消多少”连续可调，这个量必须带一个连续变化的角度。概率幅是被实验逼出来的最小升级，没有比它更省事的方案。

再看这支箭头在两次测量之间干什么。第 42 章说过，能量为 \(E\) 的状态对应频率 \(f = E/h\)。在量子力学里，这句话的精确版本是：一支长度为定值的箭头，以角速度
\[
\omega = \frac{E}{\hbar}, \qquad \hbar = \frac{h}{2\pi} \approx 1.05\times10^{-34}\,\mathrm{J\cdot s}
\]
匀速旋转。用复数乘法写出来，就是振幅随时间乘以 \(\cos\omega t + i\sin\omega t\)（这个旋转因子有个更紧凑的名字，\(e^{i\omega t}\)，欧拉公式把它和三角函数连成一体）。单个箭头的整体旋转谁都测不出来——玻恩规则只看长度；但两条路径的箭头若转得快慢不同，它们的相对相位就随时间漂移，干涉图样随之移动。能量差，原来是相位转速差。至于“是什么定律规定了各种情况下箭头怎么转”，那是一个真正的动力学方程——薛定谔方程，第 44 章的主角。本章的加法与平方规则是它的入场券。
\end{challenge}

\step{模型的边界}

新语言威力巨大，下面每一条都是写在包装上的适用范围和防伪标签。

第一，波函数是预测工具，不是物质云。\(|\psi(x)|^2\) 大，意思是“在此处探测到电子的概率大”，不是“电子的一部分摊在这里”。每一次探测永远是完整的一颗、一个点。开场照片里散开的，是成千上万次独立事件的统计分布，不是任何一颗电子的形状。

第二，概率幅不是概率。它可以为负、可以为复数，可以被另一支箭头抵消到零，它没有“多大概率发生”这一说；全宇宙没有任何仪器能直接读出某处的 \(\psi(x)\)，能读的只有计数率——而计数率正比于模的平方。把“概率幅”当成“某种更精细的概率”，是学量子力学最容易踩进去的第一个坑。

第三，量子概率不是“我们不知道”的无知概率。掌心里的硬币有确定答案，你的概率只是记账；而双缝后的电子，在“走哪条路”没有任何记录存在时，连“答案”本身都不存在——证据就是暗带：如果每颗电子都有确定的路线，概率加法就不会失效。这条边界也划出了经典规则的辖区：一旦两条路径变成原则上可区分的（比如有装置记下了路径信息），干涉项就会消失，概率加法复活，小球图像重新变得好用。这中间发生了什么，是第 45 章的主题，本章只标记边界，不展开。

第四，别把本章读成“电子在波和粒子之间来回变身”。电子从头到尾只有一种存在方式：用概率幅描述、探测时给出完整点事件。“波”和“粒子”是从宏观世界借来的两个旧词，各自只抓住了一半，交替使用它们只会制造它好像在变装的错觉。

第五，宏观世界为什么看起来那么“小球”。两个原因。其一，波长太短：一只 \(0.2\,\mathrm{kg}\) 的棒球以 \(20\,\mathrm{m/s}\) 飞行，\(\lambda = h/(mv) \approx 1.7\times10^{-34}\,\mathrm{m}\)，按 \(\Delta y = \lambda L/d\) 算出的“条纹”比原子核还密二十个数量级，任何探测手段看到的都只是平均后的光滑分布。其二，相位太脆弱：棒球表面每秒钟被无数光子、空气分子撞上，每一次碰撞都在客观上记录它的位置，把相位洗得乱七八糟，干涉项平均为零——正是数学桥层里那个“\(\cos\) 平均为零”的机制。所以“小球偷偷走固定路线”不是错误理论，而是量子理论在宏观条件下的极限近似；它的地盘极大，但边界存在，双缝实验就站在边界上。

第六，一张典型误用清单：以为“叠加”意味着宏观物体可以同时待在两处供人差遣（错，见第五条）；以为“测量之前电子其实有确定位置，只是我们不知道”（错，暗带就是这个图像的死刑判决，见模型三）；以为概率幅是“心理概率”“主观信念”（错，它是写在自然定律里的客观量，谁算都一样）；以为双缝图样是“许多电子之间的集体效应”（错，逐个发送照出条纹）。

\step{回头解释开场现象}

现在回到那台装置，把悬念一个一个关掉。

为什么探测总是一个点？因为“电子到达屏上”是一次探测事件，玻恩规则给出的是每次事件落在哪里的概率；每次事件本身是完整的——仪器要么响要么不响，没有“响半声”这一档。点，是测量的语言。

为什么点会织出条纹？每颗电子面对两条缝，只要世界上没有留下任何“它走了哪条”的记录，两条路径就各贡献一支振幅，相加、平方，得到每个位置的概率。概率的分布有峰有谷，大量独立事件按这个分布落下，条纹自动浮现。不是电子“排成”了条纹，是概率的地形有高山和深渊，落点只是依地形而行。

为什么暗带处关一条缝反而变亮？暗带处，两路振幅恰好等长反向，合成归零——不是因为那里“去不了”，而是两路贡献在那里对着干。关掉一条缝，等于撤走了一支反相的箭头，剩下那支独自给出非零概率，那个位置反而有电子落下。概率相加的世界观里这是灵异事件；振幅相加的世界观里，这是再朴素不过的向量算术。

为什么一颗一颗来、隔一分钟一颗，条纹照旧？因为干涉发生在每一颗电子自己的两路振幅之间，与任何别的电子无关。每颗电子都是独立完成整台实验的：两支箭头是它自己的两支箭头。“电子之间互相碰撞、协商、接力”的假说，从前提上就没有对象。

暗处少掉的电子去了哪里？去了亮处。亮纹中心的概率是单路的四倍而不是两倍，多出来的那份正是暗带让出的那份——全部位置的概率加起来恒等于一。还记得音箱实验里静音点上消失的声音吗？它出现在更响的地方。剧本一模一样：能量不消失，只是被相位重新分配了住址。

这套“分路—相位—复合”的规则，今天早已不是实验室里的珍奇，而是工程工具。原子干涉仪把整颗原子当作会携带相位的波包，让它沿两条路径分开再汇合：重力加速度的微小不同会让两条路径的相位差发生可测的变化，干涉条纹的移动就是重力的读数。最好的原子重力仪能测出重力加速度十亿分之一量级的变化，用来勘探地下的矿藏和空洞、给潜艇做不依赖卫星的惯性导航。本章的三条规则，就是这些仪器的全部工作原理。

也顺便给音箱实验补上完整的边界声明。它像双缝：两路贡献按相位叠加，同相更响、反相安静，挡住一路暗点复活。它不像双缝：声波的振幅是拿得到、测得准的空气压强，传播需要空气这个介质；电子的振幅测不到、不需要任何介质，它的“波”不在任何物质里起伏。类比到此为止，再往前一步就是误导。

反相抵消的思路在日常生活里还有一个熟悉的身影：降噪耳机。它用麦克风采集外界噪声，让扬声器实时发出一支反相的声波，两路叠加，到达鼓膜的压强起伏被压下去。这是经典波版本的“同相增强、反相抵消”，和你音箱实验里的静音点是同一条原理；差别仍然在于，声压是可以直接测量的实量，工程师拿示波器就能看到那两支“箭头”，而电子的振幅谁也看不到。

还有一个容易忽略的事实：这套规则对光子同样成立。把双缝实验换成单光子源来做，一颗一颗地发光子，剧本逐行重演——逐个的探测点、积累出的条纹、暗处关缝变亮。事实上，光学家常用的马赫—曾德尔干涉仪就是“光子版双缝”：一块分束镜把光分成两路，走不同的路程后再汇合，探测器读数随两路的相位差涨落。电子、光子、原子、碳六十分子，全都按同一本账记账——这也是我们说“量子态”而不说“电子态”的原因：规则是普适的，换的只是演员。

\step{脑洞实验与挑战题}

\begin{thought}[把双缝实验一路放大：从电子到你]
电子会干涉，那更大的东西呢？这不再是思想实验，而是一系列已经做了的真实实验，我们沿着尺寸往上爬。

一九九九年，蔡林格的团队让碳六十分子——六十个碳原子组成的“足球”，质量约 \(1.2\times10^{-24}\,\mathrm{kg}\)，是电子的一百多万倍——以约每秒二百米的速度飞过周期 \(100\,\mathrm{nm}\) 的光栅。它的德布罗意波长 \(\lambda = h/(mv) \approx 2.8\times10^{-12}\,\mathrm{m}\)，只有三皮米左右，比原子小几十倍。按 \(\lambda L/d\) 估算，屏上图样的特征间距在几十微米量级——实验家们真的测到了清晰的衍射图样。一颗含有六十个原子、内部还在发热振动的“小球”，照样拒绝走固定路线。

继续放大：棒球。前面算过，\(\lambda \approx 1.7\times10^{-34}\,\mathrm{m}\)。要看到它的条纹，你需要把缝距做到 \(10^{-30}\) 米量级——比原子核小一万倍还多——并且让棒球在飞行全程不与任何一个光子、任何一个空气分子交换信息。后者才是真正的死刑：棒球每一瞬间都在被环境“测量”着，路径记录无处不在，干涉项早已被平均得干干净净。

这个思想实验像什么：它像一架不断拉远的相机，让你看清“量子”与“经典”之间没有一道墙，只有渐变——波长越来越短，相位越来越留不住，小球图像越来越精确。它不像什么：它不像一张工程蓝图。你不可能靠“把棒球冷却得很小心”看到棒球的干涉，因为环境记录这条路在原理上就被堵死了；经典性是宏观物体的宿命，不是技术不够。
\end{thought}

\begin{puzzles}
\item 音箱实验里，把其中一只音箱的音量调小一半，静音点会发生什么变化？还是完全静音吗？\textbf{思路提示：}完全抵消需要两路贡献满足什么条件？画一画两支不等长、反方向的箭头，合成长度是多少？静音点会从“无声”变成什么？
\item 把双缝实验的电子加速电压提高，电子飞得更快，屏上条纹会变宽还是变窄？\textbf{思路提示：}先问动量和波长怎么变（\(\lambda = h/p\)），再代入条纹间距公式 \(\Delta y = \lambda L/d\)。再追一层：为什么历史上用电子做干涉实验总比用光难得多？
\item 把一颗玻璃球蒙上眼睛随机地从左缝或右缝弹入（每次确实只走一条，只是没人看），重复一万次，屏上为什么永远只有 \(P_1 + P_2\) 两团，而电子却有暗带？\textbf{思路提示：}“蒙上眼睛”改变的是谁知道，还是世界上存不存在记录？玻璃球每次走哪条缝这件事，在半路上有没有留下痕迹（碰撞的空气、反弹的声音、擦过的灰尘）？对玻璃球谈“两路振幅的相位差”，在操作上有意义吗？
\item 有人说：“给我足够好的仪器，我总能直接量出某处振幅的大小和相位。”试着反驳他。\textbf{思路提示：}你能制造和读取的一切信号，归根到底是什么事件？探测器的每一次响应，是“一段连续的箭头读数”还是“一次完整的计数”？计数率对应振幅的哪个属性？再想想：相位差为什么反而可测——干涉条纹的移动充当的是什么角色？
\end{puzzles}

到这里，本章交付了一台完整的预测机器：写出量子态，给每条不可区分的路径配一支振幅，相加，平方，读出概率。这台机器回答“此刻能测到什么、各有多大概率”，却对一个问题保持沉默——两次测量之间，量子态自己怎样随时间演变？那支相位箭头按什么定律转动？回答它需要一个方程，就像当年牛顿为“运动如何变化”写下他的方程一样。那是第 44 章的薛定谔方程。而测量那一刻发生的更深层的故事，则留到第 45 章。
